\SourceLeaf{152}
note in A. N. No. 110, p. 230\HistoricalNote{[*)]}{[*)}{See Vol. VI, p. 457.]}, as well —). His mistake was to leave out the mean error to be expected in the results. Moreover, by combining the different results from different combinations, in which his observations near the equator must always retain a predominant influence, he creates the appearance of a much greater precision in the final results than they actually have. This is a point Mr \textsc{Ivory} rightly criticizes. The case indeed resembles, though to a much smaller degree, the way the late \textsc{Ende} would sometimes determine polar altitudes by \textsc{Douwes}'s method. He would always combine \textls{the same} altitude near noon with \textls{different} altitudes far removed from noon and then obtain results agreeing within $0'',01$.

It is equally difficult to refute this paper by \textsc{Ivory} point by point, for it has just as little logical order. I find in it not the slightest argument against the suitability of the method of least squares. But what is one to say of the procedure that \textsc{Ivory} would substitute for it, combining the shortest pendulum length with many at high latitudes? This is quite unworthy of a sensible mathematician and is really \textsc{Ende}'s procedure\HistoricalNote{[**)]}{[**)}{The foregoing criticism refers to the following works by J. \textsc{Ivory}:
\par\hangindent1.5em 1. On the method of least squares. Phil. Mag. (Tilloch), \textsc{lxviii}, 1826, pp. 161—165.
\par\hangindent1.5em 2. On the Ellipticity of the Earth, as deduced from Experiments made with the Pendulum. Phil. Mag. (Tilloch), \textsc{lxviii}, 1826, pp. 3—10, 92—101, 241—245, 246—251, 321—326, 350—353.]}. . . . . . . .

\Needspace{6\baselineskip}
\begin{center}\textsc{Gauss} to \textsc{Encke}. Göttingen, 23 August 1831.\end{center}

. . . . . . . . . . . . . .

I was interested to learn from your letter about the course you took to justify taking the arithmetic mean. I find this course highly commendable insofar as the aim is to give an answer to the question of what to do entirely independent of any considerations of\space%
